%=====================================================================
\chapter{Unsupervised Learning and Anomaly Detection}\label{ch:unsupervised}
%=====================================================================
\locator{3}{Part III --- Machine Learning Core}

\begin{outcomes}
By the end of this chapter you will be able to:
\begin{tight}
  \item \textbf{Run} and \textbf{interpret} k-means, hierarchical and density-based clustering, and choose between them.
  \item \textbf{Select} the number of clusters using the elbow and silhouette methods.
  \item \textbf{Explain} what PCA does and read a scree plot.
  \item \textbf{Choose} an anomaly detection method appropriate to the data and the definition of ``anomalous''.
  \item \textbf{Evaluate} unsupervised results despite having no labels.
\end{tight}
\end{outcomes}

\begin{prereq}
\chref{math} (distance and scaling --- clustering is distance, so unscaled features
will decide your clusters for you) and \chref{evaluation}.
\end{prereq}

\begin{keyterms}
clustering $\bullet$ k-means $\bullet$ centroid $\bullet$ inertia $\bullet$ elbow
method $\bullet$ silhouette score $\bullet$ hierarchical clustering $\bullet$
dendrogram $\bullet$ DBSCAN $\bullet$ dimensionality reduction $\bullet$ PCA
$\bullet$ explained variance $\bullet$ t-SNE / UMAP $\bullet$ anomaly detection
$\bullet$ isolation forest $\bullet$ association rules
\end{keyterms}

\section{Learning Without Answers}

Unsupervised learning finds structure when no labels exist. That describes most
real data: nobody has labelled which of your 12{,}000 accounts are compromised, or
which students form meaningful behavioural groups. This chapter is also the direct
technical foundation of \chref{cybersecurity}, where the anomaly is the target.

\section{k-Means Clustering}

\begin{definitionbox}[k-Means]
Choose $k$. Place $k$ centroids at random. Repeat until stable: (1) assign each
point to its nearest centroid; (2) move each centroid to the mean of its assigned
points. It minimises \term{inertia} --- the total squared distance from points to
their own centroid.
\end{definitionbox}

\dsfig{%
\begin{tikzpicture}
\begin{groupplot}[
  group style={group size=4 by 1, horizontal sep=0.5cm},
  width=4.25cm, height=3.7cm,
  axis lines=left, xtick=\empty, ytick=\empty,
  title style={font=\tiny\bfseries},
  xmin=0, xmax=10, ymin=0, ymax=10,
]
\nextgroupplot[title={1. random centroids}]
\addplot[only marks, mark=*, mark size=0.9pt, gray!60] coordinates {
 (1.4,7.9)(2.2,8.6)(1.8,7.1)(2.9,8.2)(2.4,6.9)(1.1,8.4)
 (7.6,7.8)(8.4,8.5)(7.9,6.9)(8.8,7.4)(7.2,8.6)(8.2,6.4)
 (4.6,2.1)(5.4,1.4)(4.9,2.9)(5.9,2.4)(4.2,1.6)(5.6,3.1)};
\addplot[only marks, mark=x, mark size=3pt, accent, line width=1.2pt]
  coordinates {(3.2,4.2)(6.1,5.4)(5.0,7.0)};

\nextgroupplot[title={2. assign to nearest}]
\addplot[only marks, mark=*, mark size=0.9pt, secondary] coordinates {
 (1.4,7.9)(2.2,8.6)(1.8,7.1)(2.9,8.2)(2.4,6.9)(1.1,8.4)};
\addplot[only marks, mark=*, mark size=0.9pt, greenhl] coordinates {
 (7.6,7.8)(8.4,8.5)(7.9,6.9)(8.8,7.4)(7.2,8.6)(8.2,6.4)};
\addplot[only marks, mark=*, mark size=0.9pt, goldhl] coordinates {
 (4.6,2.1)(5.4,1.4)(4.9,2.9)(5.9,2.4)(4.2,1.6)(5.6,3.1)};
\addplot[only marks, mark=x, mark size=3pt, accent, line width=1.2pt]
  coordinates {(3.2,4.2)(6.1,5.4)(5.0,7.0)};

\nextgroupplot[title={3. move centroids}]
\addplot[only marks, mark=*, mark size=0.9pt, secondary] coordinates {
 (1.4,7.9)(2.2,8.6)(1.8,7.1)(2.9,8.2)(2.4,6.9)(1.1,8.4)};
\addplot[only marks, mark=*, mark size=0.9pt, greenhl] coordinates {
 (7.6,7.8)(8.4,8.5)(7.9,6.9)(8.8,7.4)(7.2,8.6)(8.2,6.4)};
\addplot[only marks, mark=*, mark size=0.9pt, goldhl] coordinates {
 (4.6,2.1)(5.4,1.4)(4.9,2.9)(5.9,2.4)(4.2,1.6)(5.6,3.1)};
\addplot[only marks, mark=x, mark size=3pt, accent, line width=1.2pt]
  coordinates {(2.1,7.9)(7.9,7.6)(5.1,2.3)};
\draw[-{Stealth[length=1.4mm]}, accent!60, dashed] (axis cs:3.2,4.2) -- (axis cs:2.1,7.9);
\draw[-{Stealth[length=1.4mm]}, accent!60, dashed] (axis cs:6.1,5.4) -- (axis cs:7.9,7.6);
\draw[-{Stealth[length=1.4mm]}, accent!60, dashed] (axis cs:5.0,7.0) -- (axis cs:5.1,2.3);

\nextgroupplot[title={4. converged}]
\addplot[only marks, mark=*, mark size=0.9pt, secondary] coordinates {
 (1.4,7.9)(2.2,8.6)(1.8,7.1)(2.9,8.2)(2.4,6.9)(1.1,8.4)};
\addplot[only marks, mark=*, mark size=0.9pt, greenhl] coordinates {
 (7.6,7.8)(8.4,8.5)(7.9,6.9)(8.8,7.4)(7.2,8.6)(8.2,6.4)};
\addplot[only marks, mark=*, mark size=0.9pt, goldhl] coordinates {
 (4.6,2.1)(5.4,1.4)(4.9,2.9)(5.9,2.4)(4.2,1.6)(5.6,3.1)};
\addplot[only marks, mark=x, mark size=3pt, accent, line width=1.2pt]
  coordinates {(2.0,7.85)(8.0,7.6)(5.1,2.25)};
\end{groupplot}
\end{tikzpicture}}%
{k-means in four steps. Assignment and update alternate until nothing moves. Because the
starting positions are random, different runs can converge differently --- which is why
\texttt{n\_init} exists and why you should fix the random seed.}{kmeans-steps}

\subsection{\texorpdfstring{Choosing $k$}{Choosing k}}

\dsfig{%
\begin{tikzpicture}
\begin{groupplot}[
  group style={group size=2 by 1, horizontal sep=1.5cm},
  width=6.6cm, height=4.4cm,
  axis lines=left, tick label style={font=\tiny}, label style={font=\tiny},
  title style={font=\scriptsize\bfseries},
]
\nextgroupplot[title={Elbow method: inertia against $k$},
  xlabel={number of clusters $k$}, ylabel={inertia},
  xmin=0.5, xmax=9.5, ymin=0, ymax=105, xtick={1,...,9}]
\addplot[primary, line width=1.3pt, mark=*, mark size=1.6pt] coordinates {
 (1,100)(2,52)(3,26)(4,21)(5,17.5)(6,15)(7,13.2)(8,11.8)(9,10.8)};
\draw[accent, dashed, line width=1pt] (axis cs:3,0) -- (axis cs:3,26);
\node[font=\tiny, accent, anchor=west] at (axis cs:3.25,40) {elbow at $k=3$};
\node[font=\tiny, text=gray!55!black, anchor=west, align=left] at (axis cs:4.6,72)
     {inertia always falls\\as $k$ rises --- at $k=n$\\it reaches zero, which\\is why you
      cannot just\\minimise it};

\nextgroupplot[title={Silhouette score against $k$},
  xlabel={number of clusters $k$}, ylabel={mean silhouette},
  xmin=0.5, xmax=9.5, ymin=0, ymax=0.85, xtick={2,...,9}]
\addplot[greenhl, line width=1.3pt, mark=*, mark size=1.6pt] coordinates {
 (2,0.52)(3,0.71)(4,0.58)(5,0.46)(6,0.41)(7,0.36)(8,0.33)(9,0.30)};
\addplot[accent, only marks, mark=*, mark size=2.4pt] coordinates {(3,0.71)};
\node[font=\tiny, accent, anchor=west] at (axis cs:3.3,0.71) {maximum at $k=3$};
\node[font=\tiny, text=gray!55!black, anchor=west, align=left] at (axis cs:5.0,0.68)
     {this one has a genuine\\maximum, so it is the\\more reliable of the two};
\end{groupplot}
\end{tikzpicture}}%
{Two ways to choose $k$. Inertia falls monotonically, so the elbow is a judgement call;
the silhouette score has a true maximum and is the better default. Both should agree with
a domain reason for the number you pick.}{choosing-k}

\begin{definitionbox}[Silhouette Score]
For each point, compare $a$ (mean distance to its own cluster) with $b$ (mean
distance to the nearest other cluster): $s = \dfrac{b-a}{\max(a,b)}$. Near $+1$
means well-clustered, near $0$ means on a boundary, negative means probably in the
wrong cluster. Average over all points.
\end{definitionbox}

\section{Hierarchical Clustering and DBSCAN}

\rowcolors{2}{white}{lightbg}
\begin{longtable}{@{}L{2.5cm}L{3.5cm}L{3.6cm}L{4.4cm}@{}}
\toprule
\rowcolor{primary!15}
& \textbf{k-means} & \textbf{Hierarchical} & \textbf{DBSCAN} \\
\midrule
\endhead
Needs $k$? & Yes, in advance & No --- cut the dendrogram anywhere & No \\
Cluster shape & Spherical only & Any & \textbf{Any} \\
Outliers & Forced into a cluster & Forced into a cluster & \textbf{Labelled as noise} \\
Scales to large $n$? & Yes & No ($O(n^2)$ or worse) & Moderately \\
Key parameter & $k$ & linkage, cut height & \texttt{eps}, \texttt{min\_samples} \\
Use when & Large data, roughly round groups & You want the hierarchy itself & Irregular shapes, or you need outliers identified \\
\bottomrule
\end{longtable}

\begin{conceptbox}[Why DBSCAN Matters for This Course]
DBSCAN defines clusters as dense regions and labels everything else \emph{noise}.
For cybersecurity and IoT that is not a side-effect --- it is the output you want.
The points k-means would have swept into the nearest cluster are precisely the
sessions or devices you were looking for.
\end{conceptbox}

\dsfig{%
\begin{tikzpicture}[font=\scriptsize]
% dendrogram
\draw[primary, line width=1pt] (0,0) -- (0,1.0) -- (1.0,1.0) -- (1.0,0);
\draw[primary, line width=1pt] (2.0,0) -- (2.0,0.7) -- (3.0,0.7) -- (3.0,0);
\draw[primary, line width=1pt] (0.5,1.0) -- (0.5,2.0) -- (2.5,2.0) -- (2.5,0.7);
\draw[primary, line width=1pt] (4.0,0) -- (4.0,1.4) -- (5.0,1.4) -- (5.0,0);
\draw[primary, line width=1pt] (1.5,2.0) -- (1.5,3.1) -- (4.5,3.1) -- (4.5,1.4);
\foreach \x/\l in {0/A, 1/B, 2/C, 3/D, 4/E, 5/F}
  \node[font=\tiny, anchor=north] at (\x,-0.05) {\l};

\draw[accent, dashed, line width=1.1pt] (-0.6,1.7) -- (5.6,1.7);
\node[font=\tiny, text=accent, anchor=west] at (5.65,1.7) {cut here $\Rightarrow$ 3 clusters};
\draw[accent!50, dashed, line width=0.9pt] (-0.6,2.6) -- (5.6,2.6);
\node[font=\tiny, text=accent!60, anchor=west] at (5.65,2.6) {cut here $\Rightarrow$ 2 clusters};

\node[font=\tiny, text=gray!55!black, rotate=90, anchor=south] at (-0.85,1.5) {distance};
\node[font=\scriptsize\bfseries, text=primary] at (2.5,3.6) {Dendrogram (hierarchical)};

% DBSCAN illustration
\begin{scope}[shift={(9.6,0)}]
  \node[font=\scriptsize\bfseries, text=primary] at (0,3.6) {DBSCAN: dense core vs.\ noise};
  \foreach \x/\y in {-1.5/0.6, -1.1/1.1, -0.7/0.5, -1.3/1.6, -0.8/1.4, -1.7/1.2,
                     -0.5/1.0, -1.2/0.2, -0.4/1.7}
    \fill[greenhl] (\x,\y) circle (1.6pt);
  \foreach \x/\y in {1.0/1.9, 1.5/2.3, 1.2/1.4, 1.8/1.7, 0.8/2.4, 1.6/1.0, 2.0/2.2}
    \fill[secondary] (\x,\y) circle (1.6pt);
  \foreach \x/\y in {0.1/0.1, 2.4/0.3, -0.2/2.9, 2.7/3.0}
    {\draw[accent, line width=1pt] (\x,\y) circle (2.4pt);}
  \node[font=\tiny, text=accent, anchor=west] at (-2.0,-0.35)
       {\faIcon{circle}~open circles = \textbf{noise} --- k-means would have hidden these};
\end{scope}
\end{tikzpicture}}%
{Left: a dendrogram, where the number of clusters is chosen after the fact by cutting at a
chosen height. Right: DBSCAN separates dense regions from noise, which makes it a
clustering algorithm and an outlier detector at once.}{dendrogram-dbscan}

\section{Dimensionality Reduction}

\begin{definitionbox}[Principal Component Analysis]
PCA finds new axes --- \term{principal components} --- that are linear combinations
of the original features, ordered so that the first captures the most variance, the
second the most of what remains, and so on. Keeping the first few gives a
lower-dimensional representation that preserves most of the structure.
\end{definitionbox}

\dsfig{%
\begin{tikzpicture}
\begin{groupplot}[
  group style={group size=2 by 1, horizontal sep=1.5cm},
  width=6.6cm, height=4.6cm,
  axis lines=left, tick label style={font=\tiny}, label style={font=\tiny},
  title style={font=\scriptsize\bfseries},
]
\nextgroupplot[title={PCA finds the directions of variance},
  xlabel={feature 1}, ylabel={feature 2},
  xmin=-4, xmax=4, ymin=-4, ymax=4, xtick=\empty, ytick=\empty]
\addplot[only marks, mark=*, mark size=1pt, primary!60] coordinates {
 (-2.8,-2.2)(-2.2,-1.9)(-1.9,-1.2)(-1.4,-1.4)(-1.0,-0.5)(-0.6,-0.9)(-0.2,0.1)
 (0.3,-0.3)(0.7,0.6)(1.1,0.4)(1.6,1.2)(2.0,1.0)(2.5,1.9)(2.9,1.6)
 (-1.6,-0.3)(1.3,-0.2)(-0.4,0.9)(0.9,1.4)};
\draw[-{Stealth[length=2.4mm]}, accent, line width=1.5pt] (axis cs:0,-0.15) -- (axis cs:3.0,2.0);
\draw[-{Stealth[length=2.4mm]}, greenhl, line width=1.3pt] (axis cs:0,-0.15) -- (axis cs:-1.0,1.35);
\node[font=\tiny, accent, anchor=west] at (axis cs:2.0,2.6) {PC1: most variance};
\node[font=\tiny, greenhl!55!black, anchor=east] at (axis cs:-1.1,1.9) {PC2};

\nextgroupplot[title={Scree plot: how many components to keep},
  xlabel={component}, ylabel={variance explained (\%)},
  xmin=0.4, xmax=8.6, ymin=0, ymax=62, xtick={1,...,8}, ybar, bar width=8pt]
\addplot[draw=primary, fill=primary!35] coordinates {
 (1,54)(2,22)(3,9)(4,5)(5,4)(6,3)(7,2)(8,1)};
\draw[accent, dashed, line width=1pt] (axis cs:2.5,0) -- (axis cs:2.5,58);
\node[font=\tiny, accent, anchor=west, align=left] at (axis cs:2.65,45)
     {keep 2 components:\\76\% of the variance};
\end{groupplot}
\end{tikzpicture}}%
{PCA. Left: the components are the directions along which the data actually varies, not
the original axes. Right: the scree plot shows where adding components stops buying much,
which is how the number to keep is chosen.}{pca}

\begin{alertbox}[PCA Requires Scaling, and Destroys Interpretability]
PCA maximises variance, so an unscaled feature measured in bytes will dominate one
measured in seconds regardless of importance --- always standardise first. And the
components are mixtures: ``PC1'' has no meaning a stakeholder can act on. Use PCA
for compression, visualisation and noise reduction, not when you must explain
which feature drove a decision.
\end{alertbox}

\begin{conceptbox}[t-SNE and UMAP: For Looking, Not for Modelling]
These produce striking 2-D pictures of high-dimensional data by preserving local
neighbourhoods. Two cautions: distances \emph{between} clusters in the picture are
not meaningful, and the layout changes with the random seed and the perplexity
setting. Use them to generate hypotheses, never as evidence, and never as features
for a downstream model.
\end{conceptbox}

\section{Anomaly Detection}

\begin{definitionbox}[Anomaly Detection]
Identifying observations that differ markedly from the majority. It is the
unsupervised counterpart to the rare-class classification of \chref{evaluation},
used when labelled examples of the rare class are unavailable --- which is the norm
for novel attacks and for equipment that has not yet failed.
\end{definitionbox}

\rowcolors{2}{white}{defbg}
\begin{longtable}{@{}L{3.2cm}L{5.0cm}L{6.0cm}@{}}
\toprule
\rowcolor{primary!15}
\textbf{Method} & \textbf{Idea} & \textbf{Best for} \\
\midrule
\endhead
Statistical ($z$, IQR) & Beyond $3\sigma$ or $1.5\times$IQR & Single variables that are roughly normal. Fails badly on long tails (\chref{descriptive}) \\
Distance / k-NN & Far from its nearest neighbours & Moderate dimensions; needs scaling \\
Density (LOF, DBSCAN) & In a sparse region relative to its neighbourhood & Clusters of differing density \\
Isolation Forest & Randomly split; anomalies get isolated in few splits & \textbf{The strong default} --- fast, scales well, high dimensions \\
Autoencoder & Train to reconstruct normal data; high reconstruction error = anomaly & Complex or sequential data (\chref{architectures}) \\
Time-series residual & Forecast the next value; large error = anomaly & Telemetry and log volumes (\chref{timeseries}) \\
\bottomrule
\end{longtable}

\begin{worked}[Evaluating Clusters Without Labels]
You cluster 12{,}000 user accounts into $k=4$ and must decide whether the result is
useful. There are no labels, so accuracy is unavailable.

\smallskip
\textbf{1. Internal validity.} Silhouette $= 0.61$. Above about 0.5 the structure is
reasonably well separated. Compare against $k=3$ and $k=5$.

\textbf{2. Stability.} Re-run on a random 80\% subsample, five times. If cluster
memberships shuffle wildly between runs, the structure is an artefact of the
sample, not of the population. This test is skipped far too often.

\textbf{3. Profile the clusters.} Compute the mean of every feature per cluster and
compare against the overall mean. Cluster 3 has 4$\times$ the failed-login rate,
2$\times$ the distinct destination ports and half the session duration.

\textbf{4. External validation --- the decisive test.} Do the clusters relate to
something you did \emph{not} cluster on? If cluster 3 also contains 70\% of the
accounts later confirmed compromised, the clustering has found something real.

\smallskip
\textbf{Verdict:} steps 1--2 say the clusters are statistically sound; steps 3--4
say whether they are \emph{useful}. A clustering can pass the first pair and fail
the second, in which case it is a mathematically valid answer to a question nobody
asked.
\end{worked}

%---------------------------------------------------------------------
\begin{fourlenses}{Finding Structure Without Labels}
\lensLA{Cluster students on engagement patterns to discover learner \emph{types} ---
steady, bursty, front-loaded, disengaging --- which no one had defined in advance
and which support different interventions. Use $k$-means with silhouette selection,
then profile each cluster in plain language. \textbf{Ethical caution:} these
clusters must never become permanent labels attached to students
(\chref{ethics}); they describe a semester's behaviour, not a person.}

\lensPM{Cluster past projects to find natural categories of work, giving better
reference classes for estimation than the official project taxonomy. With only
dozens of projects, hierarchical clustering is preferable to k-means --- the
dendrogram itself is the useful output, because a manager can see \emph{why} two
projects were grouped and overrule it.}

\lensIOT{Cluster devices by behavioural signature to find fleet sub-populations
that need different maintenance regimes. More importantly, use \emph{isolation
forests on window features} for anomaly detection: you cannot supervise a failure
mode that has never occurred, and this is the standard approach when there are no
positive labels at all. Establish each device's own baseline --- fleet-wide
thresholds miss a device that is abnormal only for itself.}

\lensSEC{This is the chapter's home domain. Unsupervised methods detect \emph{novel}
attacks that no signature covers --- precisely what supervised models trained on
past attacks cannot do. Use DBSCAN or isolation forests for user and entity
behaviour analytics, scoring each account against its own 30-day profile. Remember
the base rate from \chref{probability}: an anomaly score is not a verdict, and
tuning contamination too high will bury your analysts.}
\end{fourlenses}

%---------------------------------------------------------------------
\begin{lab}[Cluster, Reduce, Detect]
\begin{lstlisting}[language=Python]
import numpy as np, matplotlib.pyplot as plt
from sklearn.preprocessing import StandardScaler
from sklearn.cluster import KMeans, DBSCAN
from sklearn.decomposition import PCA
from sklearn.ensemble import IsolationForest
from sklearn.metrics import silhouette_score

Xs = StandardScaler().fit_transform(X)      # MANDATORY -- clustering is distance

# --- choose k by silhouette, not by eye --------------------------------
for k in range(2, 9):
    km = KMeans(n_clusters=k, n_init=10, random_state=0).fit(Xs)
    print(f"k={k}  inertia={km.inertia_:9.1f}  silhouette={silhouette_score(Xs, km.labels_):.3f}")

# --- stability check: do the clusters survive resampling? --------------
rng = np.random.default_rng(0)
for _ in range(3):
    idx = rng.choice(len(Xs), size=int(0.8*len(Xs)), replace=False)
    km2 = KMeans(n_clusters=4, n_init=10, random_state=0).fit(Xs[idx])
    print("  subsample silhouette:", round(silhouette_score(Xs[idx], km2.labels_), 3))

# --- PCA for visualisation and for the scree plot ----------------------
p = PCA().fit(Xs)
print("variance explained:", (p.explained_variance_ratio_[:5]*100).round(1))
proj = PCA(n_components=2).fit_transform(Xs)

# --- anomaly detection with no labels at all ---------------------------
iso = IsolationForest(contamination=0.01, random_state=0).fit(Xs)
score = -iso.score_samples(Xs)              # higher = more anomalous
top = np.argsort(score)[::-1][:20]
print("20 most anomalous rows:", top)

plt.scatter(proj[:,0], proj[:,1], s=6, c=score, cmap="viridis")
plt.colorbar(label="anomaly score"); plt.xlabel("PC1"); plt.ylabel("PC2"); plt.show()
\end{lstlisting}
\end{lab}

%---------------------------------------------------------------------
\begin{checkpoint}
\begin{tightnum}
  \item Why does inertia always fall as $k$ increases, and what does that imply
        about using it alone to choose $k$?
  \item You need clusters of irregular shape and want outliers identified rather
        than absorbed. Which algorithm, and why not k-means?
  \item PCA on unscaled data returns a first component that is almost entirely one
        feature. What happened?
\end{tightnum}
\end{checkpoint}

%---------------------------------------------------------------------
\begin{chaptersummary}
\begin{tight}
  \item k-means is fast and requires $k$ in advance; it assumes spherical clusters
        and forces every point into one.
  \item Choose $k$ by silhouette (which has a true maximum) supported by the elbow
        and by a domain reason.
  \item Hierarchical clustering gives a dendrogram you can cut anywhere; DBSCAN
        finds arbitrary shapes and labels noise explicitly.
  \item PCA re-expresses data along directions of maximum variance; the scree plot
        says how many components to keep. Scale first; expect to lose
        interpretability.
  \item t-SNE and UMAP are for looking, not for modelling.
  \item Isolation forests are the strong default for anomaly detection when no
        labels exist. Validate clusters by internal score, stability, profiling and
        an external variable.
\end{tight}
\end{chaptersummary}

\begin{reviewq}
  \item \textbf{(MCQ)} Which algorithm explicitly labels points as noise?
        (a)~k-means (b)~PCA (c)~DBSCAN (d)~hierarchical clustering.
  \item \textbf{(MCQ)} A silhouette score of $-0.2$ for a point means:
        (a)~it is well clustered (b)~it is on a boundary (c)~it is probably in the
        wrong cluster (d)~it is an outlier.
  \item \textbf{(Short)} Explain why PCA requires standardised features.
  \item \textbf{(Short)} Give two reasons not to use a t-SNE plot as evidence in a
        report.
  \item \textbf{(Short)} Why is unsupervised anomaly detection the right tool for
        novel attacks, when supervised classification is not?
  \item \textbf{(Applied)} You cluster 8{,}000 IoT devices and get four clusters
        with silhouette 0.44. Describe the full validation procedure you would run
        before presenting the result, including what would make you discard it.
  \item \textbf{(Applied)} A scree plot shows components explaining 41\%, 19\%,
        14\%, 6\%, 5\%, 4\%, \dots\ of variance. How many components would you keep
        and why? What would you check before using the reduced representation in a
        model that must be explained to auditors?
\end{reviewq}
